% Fragmento compartido por VII y CRISTAL; incluir despues del cuerpo Feigenbaum.
% Procedencia: resultados recuperados, con demostraciones elementales reunidas.
% Perfil, Jacobi y escalado: 00_antecedentes_operatorios.tex, 03, 06 y 07
% del presente paquete; no se cambia su selector ni su dominio uniforme.
% Constitutiva: CRISTAL_ES/main_autosuficiente.tex:58738--58840,58864--58894.
% Inversion: CRISTAL_ES/sections/cr28_restitucion_constitutiva.tex:18--59.
% Catalan: el mismo propietario, 277--373; identidad integral, unicidad y serie.
% Corte material: CUMPLIMIENTO_EDITORIAL_20260928/fuentes/CRISTAL_ES.
% No contiene una nueva calibracion ni una nueva campana de calculo.
\section{Articulación de los lectores crítico, espectral y constitutivo}
\label{hmtfg:articulacion:seccion}

La composición APP--TRIT--TPK precede a los lectores de esta sección.
APP conserva en cada célula las evaluaciones aditiva y multiplicativa,
con sus residuos y cocientes; la reducción trítica determina el régimen,
la orientación y el acarreo de la transición. La composición efectiva
de selección, transporte y actualización de memoria produce el estado
enriquecido y sus prolongaciones, desarrollados en
\ref{hmtfg:app-trit-transporte} y
\ref{hmtfg:subsubsec:tres-vueltas-ext-ch}.
La holonomía nonádica compone esas transiciones: el retorno de fase
avanza la memoria y conserva la historia. Los lectores actúan sobre la
misma estructura discreta conjunta del continuo, con sus cinco
construcciones consustanciales y con las dos publicaciones del mismo
antecedente: evaluación ordenada e incidencia de frontera.

La rueda orientada y los canales angulares son salidas de esta
composición. Sus coordenadas HMT ya generadas se utilizan en operaciones
posteriores, con sus marcas y dominios; no son cifras metrológicas
introducidas para seleccionar una respuesta. La publicación ordenada
permite evaluar las fórmulas siguientes, mientras el estado de procedencia
conserva hoja, ruta, residuo, cociente, acarreo, frontera y memoria.
Las restituciones que se prueban recuperan el objeto de su dominio
declarado; la representación evaluada permanece acompañada de esos
registros.

\subsection{La separación crítica como lectura asintótica intrínseca}
\label{hmtfg:articulacion:separacion}

En el sector uniforme de la rueda dodecafásica, los representantes
orientados y sus pesos son
\[
 \phi_m=\frac{(3m-1)\pi_{\mathrm{HMT}}}{18},\qquad
 w_m=\frac1{12},\qquad 1\le m\le12.
\]
El segundo momento se normaliza mediante
\(\sum_{m=1}^{12}(3m-1)^2=5394=6\cdot899\).
De este modo, el factor angular común se cancela y quedan
\begin{equation}
 k_m=\frac{2(3m-1)}{\sqrt{899}},\qquad
 u_0(z)=\frac1{12}\sum_{m=1}^{12}\bigl(1-\cos(k_mz)\bigr),
 \qquad f_\lambda(z)=1-\lambda u_0(z).
 \label{hmtfg:articulacion:perfil}
\end{equation}
La variable \(\lambda\) recorre la familia después de fijados el
lector, sus pesos y la normalización. En
\(0<\lambda\le2/u_0(1)\), el dominio dinámico es \([-1,1]\).
La primera raíz es \(\lambda_1=1/u_0(1)\); para \(n\ge2\),
\(\lambda_n\) es la menor raíz mayor que \(\lambda_{n-1}\)
cuyo punto crítico tiene período exacto \(2^n\).
La clasificación, la existencia de cada mínimo y la identificación
métrica de estos centros se han demostrado en
\ref{hmtfg:escalado:15}--\ref{hmtfg:escalado:17}.

\begin{proposition}[Lectura estructural de la constante de escalado]
\label{hmtfg:articulacion:delta}
Para el lector uniforme \eqref{hmtfg:articulacion:perfil} y su selector
de primeras raíces, la constante de separación crítica es
\begin{equation}
 \delta_{\mathrm{HMT,crit}}
 :=\lim_{n\to\infty}
 \frac{\lambda_n-\lambda_{n-1}}{\lambda_{n+1}-\lambda_n}
 =\delta_{\mathrm F}.
 \label{hmtfg:articulacion:limite}
\end{equation}
Es el invariante asintótico de separación de parámetros de los retornos
críticos de ese lector, con su orden de selección conservado.
\end{proposition}
\begin{proof}
El teorema de escalado \eqref{hmtfg:escalado:eq:67} proporciona
\(\lambda_*-\lambda_n=C\delta_{\mathrm F}^{-n}(1+e_n)\),
con \(C>0\) y \(e_n\to0\). Restando dos términos consecutivos,
\[
 \lambda_n-\lambda_{n-1}
 =C\delta_{\mathrm F}^{-n}
   \bigl[\delta_{\mathrm F}-1+
          \delta_{\mathrm F}e_{n-1}-e_n\bigr].
\]
El cociente de esta igualdad y la del índice \(n+1\) converge a
\(\delta_{\mathrm F}\). La identidad eventual de los centros locales
con las primeras raíces, demostrada en
\eqref{hmtfg:escalado:eq:66}, permite aplicar el escalado precisamente
a la sucesión seleccionada y no a otra cascada elegida después.
\end{proof}

Esta formulación conserva más información que el numeral terminal:
identifica la rueda, su lectura, la familia y el selector cuyo límite
se evalúa. Su interpretación es intrínseca a esa composición HMT;
el reconocimiento de la constante universal utiliza después el teorema
analítico aplicado al lector ya construido.

\subsection{La medida finita y la restitución espectral del perfil}
\label{hmtfg:articulacion:jacobi}

Las frecuencias cuadradas producen la probabilidad atómica
\begin{equation}
 s_m=k_m^2=\frac{4(3m-1)^2}{899},\qquad
 \nu_0=\frac1{12}\sum_{m=1}^{12}\delta_{s_m},\qquad
 M_j=\int s^j\,d\nu_0(s).
 \label{hmtfg:articulacion:medida-finita}
\end{equation}
Aquí \(\delta_{s_m}\) es la masa de Dirac en un nodo, distinta
de la constante \(\delta_{\mathrm F}\). La medida conserva las
doce frecuencias y sus pesos. Para un polinomio no nulo de grado menor
que \(r\le12\),
\[
 \int p(s)^2\,d\nu_0(s)
 =\frac1{12}\sum_{m=1}^{12}p(s_m)^2>0,
\]
pues un polinomio de grado menor que doce no puede anularse en los
doce nodos distintos. Así, las matrices de Hankel
\((M_{i+j})_{0\le i,j<r}\) son positivas definidas.
El polinomio \(\prod_m(s-s_m)\) tiene norma nula; el espacio
\(L^2(\nu_0)\) tiene dimensión doce.

\begin{proposition}[Representación exacta del lector crítico]
\label{hmtfg:articulacion:representacion}
Sean \(\psi_0,\ldots,\psi_{11}\) los polinomios ortonormales
de coeficiente principal positivo para \(\nu_0\), con
\(\psi_0=1\). La multiplicación por \(s\) tiene matriz de
Jacobi \(J_{12}\), real simétrica y de subdiagonal positiva.
Con \(e_0=(1,0,\ldots,0)^{\mathsf T}\),
\begin{equation}
 u_0(z)=\int\bigl(1-\cos(z\sqrt{s})\bigr)\,d\nu_0(s)
 =e_0^{\mathsf T}\bigl[I-\cos(z\sqrt{J_{12}})\bigr]e_0.
 \label{hmtfg:articulacion:perfil-jacobi}
\end{equation}
\end{proposition}
\begin{proof}
La positividad anterior permite ortogonalizar hasta grado once.
Si \(i<j-1\), la autoadjunción de la multiplicación da
\(\langle s\psi_j,\psi_i\rangle
=\langle\psi_j,s\psi_i\rangle=0\).
Por simetría, sólo quedan la diagonal y sus dos vecinas; el cociente
positivo de los coeficientes principales fija la positividad de la
subdiagonal. Definamos
\[
 Q_{mj}=\frac{\psi_j(s_m)}{\sqrt{12}},\qquad
 D=\operatorname{diag}(s_1,\ldots,s_{12}).
\]
La ortonormalidad da \(Q^{\mathsf T}Q=I\), y evaluar la
recurrencia en los nodos da \(DQ=QJ_{12}\). Como \(Q\) es
cuadrada, \(J_{12}=Q^{\mathsf T}DQ\). Sus autovalores son
positivos, por lo que la raíz positiva y el coseno se calculan en
esta diagonalización. Finalmente,
\(Qe_0=(1,\ldots,1)^{\mathsf T}/\sqrt{12}\), lo que prueba
\eqref{hmtfg:articulacion:perfil-jacobi}.
\end{proof}

La realización por \((J_{12},e_0)\) recupera el mismo perfil;
por ello, al conservar la familia y el selector, recupera también
sus retornos y sus primeras raíces. La positividad espectral
justifica esta representación. La transversalidad paramétrica
utiliza además la derivada del retorno y sus productos orbitales,
desarrollados en \ref{hmtfg:escalado:14.3} y
\ref{hmtfg:escalado:6}. Son propiedades de operaciones diferentes:
la demostración del escalado conserva ambas y no sustituye el
control de la tangente por la positividad de los momentos.

\subsection{Los canales constitutivos y su inversión etiquetada}
\label{hmtfg:articulacion:constitutiva}

El lector constitutivo actúa sobre otro objeto tipado del mismo
soporte: el transporte angular de dos hojas. Sus coordenadas ya
generadas son \(x=A_{\rm rad}\) y \(y=C^*_{\rm rad}\), en la
cámara contractiva \(x>|y|\). La involución de hojas
\(R=R^*=R^{-1}\) determina los proyectores marcados
\(P_\pm=(I\pm R)/2\). La representación del transporte es
\begin{equation}
 q_+=\exp[-(x+y)],\qquad q_-=\exp[-(x-y)],\qquad
 T=q_+P_++q_-P_-,\qquad 0<q_\pm<1.
 \label{hmtfg:articulacion:canales}
\end{equation}
La evaluación exponencial realiza esos canales después de su
construcción angular. Las incidencias \(90\) y \(120\) del
transporte fijan la respuesta
\[
 \mathcal V=(I-T^{90})(I-T^{120})^{-1}.
\]
El inverso existe porque \(1-q_\pm^{120}>0\).
Escribiendo \(s=q^{30}\), con \(30=\gcd(90,120)\), se obtiene
\begin{equation}
 f(s)=\frac{1-s^3}{1-s^4}
     =\frac{1+s+s^2}{1+s+s^2+s^3},\qquad
 r_\pm=f(q_\pm^{30}),\qquad
 \mathcal V=r_+P_++r_-P_-.
 \label{hmtfg:articulacion:respuesta}
\end{equation}
En esta subsección \(f\) designa la respuesta constitutiva;
\(f_\lambda\) sigue designando el mapa dinámico de
\eqref{hmtfg:articulacion:perfil}. Sus argumentos y dominios
son distintos.

\begin{proposition}[Restitución de los canales y de la orientación]
\label{hmtfg:articulacion:inversion}
La función constitutiva \(f:(0,1)\to(3/4,1)\) es una biyección
decreciente. Su inversa \(\psi\) asigna a \(r\) la única raíz
\(s\in(0,1)\) de
\begin{equation}
 r s^3+(r-1)(1+s+s^2)=0.
 \label{hmtfg:articulacion:cubica}
\end{equation}
Las lecturas constitutivas normalizadas, con sus hojas etiquetadas,
\[
 \widehat\varepsilon=r_+^2,\qquad
 \widehat\mu=r_-^2,\qquad
 \widehat Z=\frac{r_-}{r_+},\qquad
 \widehat c=\frac1{r_+r_-},
\]
con \(\widehat\varepsilon,\widehat\mu\in(9/16,1)\), permiten recuperar
\begin{equation}
 \begin{aligned}
 r_+&=\sqrt{\widehat\varepsilon},&
 r_-&=\sqrt{\widehat\mu},\\
 s_\pm&=\psi(r_\pm),& q_\pm&=s_\pm^{1/30}.
 \end{aligned}
 \label{hmtfg:articulacion:recuperacion-canales}
\end{equation}
\begin{equation}
 x=-\frac12\log(q_+q_-),\qquad
 y=\frac12\log\frac{q_-}{q_+}.
 \label{hmtfg:articulacion:recuperacion-angular}
\end{equation}
\end{proposition}
\begin{proof}
La derivación del cociente racional da
\begin{equation}
 f'(s)=-\frac{s^2(3+2s+s^2)}{(1+s+s^2+s^3)^2}<0
 \qquad(0<s<1).
 \label{hmtfg:articulacion:monotonia}
\end{equation}
Los límites \(f(0^+)=1\) y \(f(1^-)=3/4\), junto con la
continuidad, prueban la biyección. Multiplicar \(f(s)=r\) por
el denominador positivo produce \eqref{hmtfg:articulacion:cubica}.
La positividad de \(r_\pm\) determina las raíces cuadradas
y la de \(s_\pm\) determina las raíces de orden treinta.
Por \eqref{hmtfg:articulacion:canales},
\(\log q_+=-x-y\) y \(\log q_-=-x+y\);
sumar y restar recupera \(x,y\).
Recíprocamente, las raíces recuperadas pertenecen a \((0,1)\),
de modo que \(x\pm y=-\log q_\pm>0\); la pareja obtenida
permanece en la cámara \(x>|y|\).
\end{proof}

En particular,
\(\widehat\mu\widehat\varepsilon=\widehat c^{-2}\) y
\(\widehat\mu/\widehat\varepsilon=\widehat Z^2\).
El intercambio de las etiquetas permuta \(q_+\) y \(q_-\),
conserva \(x\) y cambia el signo de \(y\). La restitución
requiere esa orientación: el producto aislado \(r_+r_-\)
conserva la propagación común, pero no determina las dos respuestas
etiquetadas. Las raíces positivas y la inversión cúbica recuperan
los canales dentro del lector fijado; el registro completo de su
generación permanece en el estado enriquecido.

\subsection{El momento de Catalán y la restitución funcional}
\label{hmtfg:articulacion:catalan}

Las mismas incidencias determinan una relación entre funciones,
antes de evaluar cualquiera de los canales. Para \(0<s<1\), sea
\(\mathcal D=s\,d/ds\), y definamos el momento orientado
\begin{equation}
 \mathcal C_2(s)
 =\int_0^s\log\frac{s}{t}\,
     \frac{1+t}{1+t+t^2}f(t)\,dt
 =\int_0^s\frac{\log(s/t)}{1+t^2}\,dt.
 \label{hmtfg:articulacion:integral}
\end{equation}
La segunda igualdad procede de la factorización
\(1+t+t^2+t^3=(1+t)(1+t^2)\). El peso continuo y positivo
\(dt/(1+t^2)\) es así producido por la composición racional
del lector constitutivo, y no por una identificación con
\(\nu_0\).

\begin{proposition}[Inversión diferencial, unicidad y lectura extrema]
\label{hmtfg:articulacion:restitucion-catalan}
El momento \(\mathcal C_2\) es la única función
\(F\in C^2((0,1))\) que satisface
\begin{equation}
 \mathcal D^2F(s)=\frac{s(1+s)}{1+s+s^2}f(s),\qquad
 \lim_{s\downarrow0}F(s)=0.
 \label{hmtfg:articulacion:problema-diferencial}
\end{equation}
La respuesta constitutiva se restituye por
\begin{equation}
 \mathcal D^2\mathcal C_2(s)=\frac{s}{1+s^2},\qquad
 f(s)=\frac{1+s+s^2}{s(1+s)}\,
          \mathcal D^2\mathcal C_2(s).
 \label{hmtfg:articulacion:inversa-diferencial}
\end{equation}
Su extensión continua al intervalo cerrado tiene la serie
\begin{equation}
 \mathcal C_2(s)=\sum_{j=0}^{\infty}
       \frac{(-1)^j s^{2j+1}}{(2j+1)^2},\qquad 0\le s\le1,
 \qquad
 \mathcal C_2(1)=\sum_{j=0}^{\infty}\frac{(-1)^j}{(2j+1)^2}
 =G_{\rm Cat}.
 \label{hmtfg:articulacion:serie-catalan}
\end{equation}
\end{proposition}
\begin{proof}
El integrando de \eqref{hmtfg:articulacion:integral} es no negativo
y
\[
 0\le\mathcal C_2(s)\le\int_0^s\log(s/t)\,dt=s.
\]
Esto prueba integrabilidad en cero y el límite requerido.
La identidad \(\log(s/t)=\int_t^s du/u\), seguida del
intercambio de integrales no negativas, proporciona
\[
 \mathcal C_2(s)
 =\int_0^s\frac1u\left(\int_0^u\frac{dt}{1+t^2}\right)du.
\]
En todo intervalo compacto del dominio abierto se aplica el
teorema fundamental del cálculo. Por tanto,
\[
 \mathcal D\mathcal C_2(s)=\int_0^s\frac{dt}{1+t^2},
 \qquad
 \mathcal D^2\mathcal C_2(s)=\frac{s}{1+s^2}.
\]
La factorización racional usada en
\eqref{hmtfg:articulacion:integral} da
\eqref{hmtfg:articulacion:problema-diferencial}; despejar \(f\)
da \eqref{hmtfg:articulacion:inversa-diferencial}, cuyos
denominadores son positivos en \((0,1)\).

Si dos soluciones tienen el límite nulo requerido, su diferencia
\(H\) cumple \(\mathcal D^2H=0\). El cambio \(z=\log s\)
transforma esta ecuación en \(d^2H(e^z)/dz^2=0\), de donde
\(H(s)=a\log s+b\). La existencia del límite nulo obliga a
\(a=b=0\), y prueba la unicidad.

Para \(s<1\), la expansión geométrica de \((1+t^2)^{-1}\)
puede integrarse término a término con el peso logarítmico, pues
\[
 \int_0^s t^{2j}\log(s/t)\,dt
 =\frac{s^{2j+1}}{(2j+1)^2},\qquad
 \sum_{j\ge0}\frac{s^{2j+1}}{(2j+1)^2}<\infty.
\]
Esto prueba la serie en el interior. La cota
\(1/(2j+1)^2\) da convergencia absoluta y uniforme de esa
serie en \([0,1]\). En la integral, el integrando extendido
por cero para \(t>s\) está dominado por \(\log(1/t)\)
en \((0,1)\). La convergencia dominada permite alcanzar
\(s=1\) y coincide con la extensión de la serie. Su valor
extremo es precisamente la serie de Catalán indicada.
\end{proof}

La inversión diferencial utiliza la familia \(\mathcal C_2(s)\),
no solamente su valor extremo \(G_{\rm Cat}\). En particular,
\(\mathcal D^2\mathcal C_2(s)\to1/2\) cuando \(s\uparrow1\).
La serie diferenciada dos veces con \(\mathcal D\) tiene allí
un límite radial de Abel; la serie del propio momento es
absolutamente convergente en la frontera. Esta distinción mantiene
las condiciones precisas de cada paso al límite.

\subsection{Composición de las restituciones y alcance conjunto}
\label{hmtfg:articulacion:conclusion}

Las interrelaciones anteriores tienen operaciones y dominios
explícitos. Los momentos de la medida atómica \(\nu_0\)
producen \(J_{12}\), y su cálculo funcional recupera el perfil
\(u_0\). El perfil y el selector de primeras raíces determinan
la sucesión cuyo escalado publica \(\delta_{\mathrm F}\).
Por otra parte, la función constitutiva \(f\) produce
\(\mathcal C_2\) mediante el momento integral, y
\eqref{hmtfg:articulacion:inversa-diferencial} recupera \(f\)
a partir de esa familia. Sus evaluaciones etiquetadas en
\(s_\pm=q_\pm^{30}\) determinan las respuestas constitutivas;
la inversión cúbica recupera ambos canales y sus coordenadas
angulares. La lectura extrema del momento da \(G_{\rm Cat}\).

La medida \(\nu_0\) es positiva, atómica y de rango doce.
El momento de Catalán integra un peso positivo continuo y posee
una expansión alternada. Los signos de esa expansión son
coeficientes de la serie geométrica; no son pesos negativos de
\(\nu_0\), ni convierten la integral positiva en la misma
representación espectral. Cada lector conserva su soporte,
orientación, normalización y operación de evaluación.

La unidad reside en la estructura APP--TRIT--TPK que genera y
transporta sus antecedentes; las relaciones funcionales se
establecen por las composiciones demostradas. La igualdad de
origen no identifica operadores de dominios diferentes.
En particular, esta articulación no postula una fórmula escalar
que determine \(\delta_{\mathrm F}\) a partir de Catalán y de
otras constantes físicas. Sitúa el escalado crítico junto a
restituciones efectivas entre lectores del mismo soporte,
conservando la información requerida por cada dirección.
